
\subsubsection{Key Findings}

The 64.7\% mediation proportion suggests preprocessing entropy is a primary pathway through which metadata completeness relates to variance reduction. This has practical implications: efforts to improve reproducibility may benefit from prioritizing preprocessing specification in dataset documentation.

Unlike existing reproducibility tools that assess experiments post-hoc, the approach examined here enables prediction before code is written.

\subsubsection{Limitations}

\textbf{Synthetic Data}: All results derive from synthetic data due to OpenML API timeout. The synthetic data follows expected distribution patterns, but results may not generalize to real OpenML data. Real-world replication is the priority next step.

\textbf{Observational Design}: Metadata completeness cannot be randomized, so causal claims are not supported. The language throughout uses ''associated with'' and ''predicts'' rather than ''causes.''

\textbf{OpenML-Specific Scope}: Generalization to HuggingFace, UCI, or deep learning benchmarks is untested.

\textbf{Model Convergence}: The mixed-effects model showed convergence warnings, which is common with complex random effects structures but may affect estimate precision.

\textbf{Sub-prediction P2b Failure}: Hyperparameter entropy varied by metadata quartile contrary to prediction. Interpretation of mechanism results should account for this unexpected pattern.

\subsubsection{Broader Impact}

\textbf{Potential Benefits}: If validated with real data, the approach could enable informed dataset selection, provide evidence-based documentation guidance, and shift reproducibility from reactive to proactive.

\textbf{Potential Risks}: Metric gaming (superficial documentation to optimize scores); selection bias (avoiding important but poorly-documented domains); misinterpretation of correlational results as causal.

